% SPDX-License-Identifier: CC-BY-SA-4.0
% Author: Matthieu Perrin

\begingroup

\begin{frame}{Résumé}

  \begin{block}{Théorème}
    \vspace{-3mm}
    \begin{center}
      Il n'existe pas d'implémentation lock-free linéarisable de la séquence\\
      qui n'utilise que des lectures et écritures sur des variables partagées.
    \end{center}
    \vspace{-2mm}
  \end{block}

  \begin{exampleblock}{Démonstration}
    \begin{itemize}
    \item[Thm1 :] Si on pouvait implémenter une séquence, \\ \hspace{5mm} alors on pourrait résoudre le consensus {\color{exampleColor}entre deux threads}
    \item[Thm2 :] Or, on ne peut pas résoudre le consensus entre deux threads
      \begin{description}
      \item[Introduction :] Supposons qu'il existe un tel algorithme $A$ ;
      \item[Lemme 1 :] $A$ possède une configuration critique ;
      \item[Lemme 2 :] $A$ ne possède pas de configuration critique ;
      \item[Conclusion :] absurde : $A$ n'existe pas !
      \end{description}
    \item[Conc :] Donc on ne peut pas implémenter une séquence
    \end{itemize}
  \end{exampleblock}

  \vspace{-2mm}
  \begin{alertblock}{Solution}
    \vspace{-1mm}
    \begin{itemize}
    \item Il faut ajouter de la puissance de calcul au modèle multi-threads
    \item Toutes les instructions spéciales se valent-elles ? 
    \end{itemize}
  \end{alertblock}

\end{frame}

\endgroup
\endinput
